\section{Introduction}
\label{sec:introduction}

We study when a map into a product of cubes can be moved onto many
coordinate boundaries while retaining every boundary value already
present.  The input spaces may have arbitrarily large, or infinite,
covering dimension.  The obstruction we use is a relative point class
in complex cobordism, tested against a family of vectors supported in
only a few cube factors.  Our main results convert the vanishing of
this class into compression on arbitrarily large Cartesian powers.

The proof has two coherence requirements.  A cobordism-null map on
one subcomplex must give a single stable nullhomotopy on a union of
products, including all their intersections.  Subsequently, stable
avoidance data on small coordinate boxes must agree on adjacent
simplices.  We prove both compatibility statements explicitly.  A
padding construction then removes the dimension restriction needed
to turn stable avoidance into an actual map.

\subsection{Coherent products}

Write $\mathbb S$ for the sphere spectrum and $MU$ for the complex
cobordism spectrum.  A based map is \emph{stably null} if its suspension
spectrum is nullhomotopic.  The unit $\mathbb S\to MU$ sends a stable
cohomotopy class to its complex-cobordism class.  Our first result
concerns the difference between vanishing on a subcomplex and
vanishing on an entire family of products.

\begin{theorem}[Filtered products]
\label{thm:filtered-product}
Let $G\subset E$ be a finite pointed CW pair, let $d\geq0$ be an integer,
and let $v:E\to S^d$ be a based map.  Suppose that the composite
\[
\Sigma^\infty G\longrightarrow\Sigma^\infty E
 \xrightarrow{\Sigma^\infty v}\Sigma^d\mathbb S
 \longrightarrow\Sigma^d MU
\]
is nullhomotopic.  For $\ell\geq1$, let $U_\ell\subset E^{\wedge\ell}$
be the union of the subcomplexes in which at least
$\lceil\ell/2\rceil$ factors are required to belong to $G$.
Then, for every sufficiently large $\ell$, the restriction
\[
v^{\wedge\ell}|_{U_\ell}:U_\ell\longrightarrow S^{d\ell}
\]
is stably nullhomotopic.
\end{theorem}

The conclusion concerns the union itself.  It therefore supplies a
nullhomotopy that can be pulled back along any map into that union.
The fraction $1/2$ may be replaced by any fixed positive fraction;
see \Cref{rem:filtered-product-proportions}.  The proof also gives a
more general compatible lift for an arbitrary commutative ring
spectrum, before invoking a vanishing theorem for $MU$.  Its target is
the total fiber of a cube of unit maps.  To produce one null cone over
the punctured cube, we use simplex parameters that collapse the entire
constrained union at once; the number of restricted factors determines
how many simplex degrees collapse.

\subsection{Cube slots and the torus criterion}

For positive integers $t,n$, put
\[
 C(t,n)=([-1,1]^t)^n,\qquad T(t,n)=([-2,2]^t)^n.
\]
A point has $n$ ordered \emph{slots}, each a vector in $\mathbb R^t$.
A boundary slot belongs to the whole boundary $\partial[-1,1]^t$.
Preserving it means preserving its entire vector, including scalar
coordinates that lie in $(-1,1)$.  All distances between cube points
use the maximum norm.  For a real number $s$, define the sparse
parameter arrangement
\[
 Z_s(t,n)=\{z\in T(t,n):\#\{j:z_j\ne0\}\le s\}.
\]
Thus it permits at most $\lfloor s\rfloor$ nonzero whole slots;
$Z_s=\varnothing$ for $s<0$, and $Z_s=T(t,n)$ for $s\ge n$.
For $0\le s<n$ its dimension is $t\lfloor s\rfloor$.
For $f:Y\to C(t,n)$, the map $f^{\times u}:Y^u\to C(t,un)$
concatenates the values of $f$, in copy order and then slot order.

The following torus data will turn a bundle embedding into a vanishing
point class.  Fix $c=1/32$.  For each positive integer $p$, choose a
smooth embedding
\[
 L\cong(\mathbb T^2)^p\subset(-c,c)^{2p+1}
\]
with a framed real normal line.  Choose line bundles $\lambda_i$, each
pulled back from the $i$th two-torus factor and embedded in
$\mathbb C^2$, so that
\[
 H=\bigoplus_{i=1}^p\lambda_i\subset L\times\mathbb C^{2p},
 \qquad
 \left\langle\prod_{i=1}^p c_1(\lambda_i),[L]\right\rangle=1.
\]
\Cref{lem:framed-torus} constructs these data.  Here $p$ is a rank
parameter, with no primality assumption.

\begin{notation}[Torus locus and slot bundle]
\label{not:torus-bundle}
On $L^n$ put
\[
 H_n=\bigoplus_{j=1}^n\operatorname{pr}_j^*H
       \subset L^n\times\mathbb C^{2pn}.
\]
For $f:Y\to C(2p+1,n)$ and $Z\subset T(2p+1,n)$, let
\[
 \mathcal T(f,Z)=\{(y,z)\in Y\times Z:f(y)-z\in L^n\},
 \qquad d_f(y,z)=f(y)-z.
\]
The \emph{slot bundle} on this locus is $E_f=d_f^*H_n$, of complex
rank $pn$, with its fixed embedding in the trivial bundle of rank
$2pn$.
\end{notation}

\Needspace{18\baselineskip}
\begin{theorem}[Torus compression]
\label{thm:torus-compression}
Let $p,n\ge1$ be integers, put $t=2p+1$, and choose $L,H$ as above.
Let $0<a<1/16$, let $Y$ be any compact metrizable space, and let
$f:Y\to C(t,n)$ be continuous.  Suppose there is one continuous
complex bundle monomorphism
\[
 E_f\longrightarrow\mathcal T(f,Z_{an}(t,n))\times\mathbb C^K
\]
over the entire locus $\mathcal T(f,Z_{an}(t,n))$, for a nonnegative
integer $K<2pn$.  Then for every integer $U\ge1$ there are an integer
$u\ge U$ and a continuous map
\[
 G:Y^u\longrightarrow C(t,un)
\]
such that at least $(a/8)un$ slots of $G(\mathbf y)$ lie in
$\partial[-1,1]^t$ for every $\mathbf y=(y_0,\ldots,y_{u-1})$.
Moreover, for $0\le b<u$ and $1\le j\le n$,
\[
 f_j(y_b)\in\partial[-1,1]^t
 \quad\Longrightarrow\quad G_{b,j}(\mathbf y)=f_j(y_b).
\]
\end{theorem}

The bundle hypothesis includes continuous dependence on the sparse
parameter as its support changes.  Since the boundary count is an
integer, its stated lower bound means at least
$\lceil(a/8)un\rceil$ slots.  The theorem also includes an empty
torus locus; on a nonempty locus the embedding itself forces $K\ge pn$.
The power $u$ can be arbitrarily large, as stated, and the theorem
requires no dimension bound on $Y$.
The strict rank inequality is essential for this general criterion:
already for $n=1$ and $K=2p$, the identity map of a cube admits the canonical
bundle embedding but no boundary-preserving compression with a positive
boundary count; see \Cref{rem:rank-threshold-sharp}.  This example does
not address the optimal boundary fraction under the strict inequality.

\subsection{Background and the proof mechanism}

The nilpotence theorem of Devinatz--Hopkins--Smith
\cite{DevinatzHopkinsSmith1988} provides the stable-homotopy background
for detecting smash-nilpotence by complex cobordism.  Here we need the
quantitative connective form developed by Mathew
\cite[Section~3.4]{Mathew2018} and recorded for $BP$ in Burklund's
appendix \cite[Appendix~B]{BurklundHahnSenger2023}.  Its cutoff is
sublinear in the degree and uniform over all connective spectra at a
fixed prime.  This uniform statement, together with explicit upper
cell bounds, converts a linearly growing Adams filtration into
vanishing.  The relation between unit-fiber powers, Adams towers and
the cobar construction is standard
\cite[Section~2.1, Propositions~2.11 and~2.14]{MathewNaumannNoel2017}.
The compatible null cone on the whole constrained product union is
constructed in \Cref{sec:filtered-products}.

The geometric step belongs to the collapse-and-duality framework of
Thom and Atiyah \cite{Thom1954,Atiyah1961Thom}.  We prove the precise
local duality needed for finite unions of coordinate planes, including
its variance under changes of boxes and arrangements.  A single
relative nullhomotopy then gives a compatible stable lift over all
simplices.  Related stable obstruction criteria for sections in a
connectivity range appear in
\cite[Theorems~4.2--4.3]{Larmore1972} and
\cite[Corollary~3.5]{KleinWilliams2007}.  Here relative support and
nested boxes keep the lift compatible and control its displacement.
Freudenthal stabilization, in the form recalled in
\cite[Corollary~4.24]{Hatcher2002}, turns it into a small actual
displacement when the parameter dimension is in the indicated range.
The singular arrangement and diagram compatibility assertions used
here are proved in \Cref{sec:stable-avoidance}.

The bundle criterion is related to the complementary-bundle Chern
obstructions used by Hirshberg--Phillips in their lower bounds for the
radius of comparison; see the proof of Theorem~3.3 in their author
preprint \cite{HirshbergPhillips2022}.  We work
with a supported point class in $MU$.  The characteristic-class
framework is classical; see Quillen \cite{Quillen1971Cobordism} and
the precise orientation and Thom formulas in
\cite[Lectures~5--6]{Lurie2010}.  If $r=pn$, the complementary bundle
has rank less than $r$, which forces the product of the $r$ square-zero
Chern classes to vanish.  A framed real Thom class and relative
excision then give vanishing of the supported point class.  These
steps retain the whole sparse parameter locus.

The resulting theorem can be used in cover constructions.  Fix a
homeomorphism from the cube to a simplex carrying boundary to boundary.
A boundary output then has at least one zero simplex coordinate.
Preserving an original boundary vector retains the entire corresponding
simplex point, and hence every coordinate that was already zero; this
is the agreement used in the companion's reconstruction lemma
\cite[Lemma~4.3]{OpenAIRadiusMeanDimension2026}.
The companion \cite[Theorem~1.1]{OpenAIRadiusMeanDimension2026}
uses \Cref{thm:torus-compression}, together with dynamical comparison
and cover reconstruction, to prove
$\rc(C(X)\rtimes_h\Z)=\tfrac12\mdim(X,h)$ for every minimal
homeomorphism $h$ of an infinite compact metrizable space $X$,
with both sides allowed to be infinite.
Here $\rc$ denotes radius of comparison and $\mdim$ denotes mean dimension.

Here is the sequence of constructions in the proof.
\begin{enumerate}
\item A point test for $f:P\to C(t,n)$ collapses the difference
      $f(x)-z$ near the origin, with $z$ in a prescribed parameter
      arrangement.  Product point tests are smash products.
      \Cref{thm:filtered-product} therefore promotes $MU$-vanishing
      of a sparse test to stable vanishing for a large batch of copies.
\item Compatible local duality converts stable test vanishing into
      a map avoiding the sparse arrangement.  An elementary radial
      normalization creates boundary slots and restores all original
      boundary vectors exactly.
\item At each new batch, we keep only the old output as a parameter.
      Its dimension is bounded by its ambient cube.  A fixed number
      of initial unchanged batches then keeps every step within the
      stable avoidance range.  The resulting boundary fraction is
      at least $a/8$ for arbitrarily large powers.
\item The torus-bundle embedding forces the $MU$ point test to vanish.
      For an arbitrary compact metrizable $Y$, finitely many auxiliary
      coordinates give a polyhedral model, while including every
      coordinate of $f$ preserves $f$ exactly.  Compression on that
      model pulls back with the same boundary agreement.
\end{enumerate}

\Cref{sec:filtered-products} proves the filtered-product theorem.
\Cref{sec:stable-avoidance} develops the point tests, local duality and
stable avoidance.  \Cref{sec:compression-cobordism} proves the general
$MU$-null boundary-compression theorem by padding.
\Cref{sec:torus-criterion} establishes the bundle criterion and the
exact compact-to-polyhedral reduction, completing
\Cref{thm:torus-compression}.

\subsection{Conventions}

Finite polyhedra and their pairs may be subdivided as needed.
A quotient $K/L$ means $K_+/L_+$, where $+$ adjoins a disjoint
basepoint; in particular, $K/\varnothing=K_+$.
Suspension spectra, smash products, cofibers and limits of spectra are
derived.  The notation $F(K,E)$ denotes the function spectrum.
We retain suspension shifts whenever connectivity or cell dimensions
are used.  A spectrum is $b$-connective when its homotopy groups in
degrees below $b$ vanish.
